\begin{papercard}{li2022}
\cf{Stages} \stagetag{A} (whole-store population rather than a sample), \stagetag{D} (repeated automated capture of store metadata and privacy labels), \stagetag{H} (reliability of self-reported labels as a harm source). A 7-page CHI extended abstract: the first measurement study of Apple's App Store privacy ``nutrition labels'', based on weekly snapshots of all U.S. apps. It is a descriptive, exploratory analysis with three research questions and no inferential statistics.

\cf{Research questions} RQ1 How promptly do developers react to the call to create a privacy label? RQ2 How often do they update the label after the initial version? RQ3 How do apps collect and use sensitive data according to their labels? \pc{li2022}{1}{RQ1 How promptly do developers react to the call of creating a privacy label?} Background: since 8 December 2020 Apple has required a label for all new apps and updates; the label is self-reported by developers, so ``the responsibility of keeping privacy labels accurate largely falls on developers'' \pc{li2022}{1}{The responsibility of keeping privacy labels accurate largely falls on developers since they are supposed to self-report their apps’ privacy practices}. A companion interview/observation study (Li et al.\ 2022, twelve iOS developers) found recurring errors from knowledge blind spots and the overhead and ambiguity of the process, which is why this paper does not test label accuracy.

\cf{Design choice: why no program analysis} Static or dynamic analysis could in principle compare labels with app behaviour, but Apple's definitions of certain terms make such a comparison impossible without access to the backend data storage, so the authors ``chose to not involve program analysis in this work'' \pc{li2022}{2}{we chose to not involve program analysis in this work}. Labels are machine-readable and standardised, which makes large-scale analysis much easier than for free-text policies or permission rationales \pc{li2022}{2}{We can directly obtain a privacy label in a machine-readable format, making the analysis a lot easier}.

\cf{Data collection} A weekly crawl of the U.S. App Store from 2 April 2021: every Thursday update the list of app IDs (removals and additions), every Friday query the privacy label and metadata for all apps in the updated list, then run error-checking scripts to find empty or incomplete records and re-crawl them \pc{li2022}{2}{Every Thursday, we start the data collection process by updating the app list to account for the removal and addition of apps}. Each label instance holds the app ID and the privacy details (data types collected; linked to the user; used to track the user). Analysis window: 2 April to 5 November 2021, 32 weeks, 1,437,605 apps \pc{li2022}{3}{based on data for 1,437,605 apps collected from April 2 to November 5, 2021}. The crawl continues beyond the paper.

\cf{Preprocessing: cohorts and ratios} \emph{New vs.\ old apps}: published after 8 Dec 2020 (new, label mandatory; 274,857 = 19.1\,\%) or before (old, forced to add one on their next update, may add voluntarily; 1,162,748 = 80.9\,\%) \pc{li2022}{3}{All new apps published on the App Store after Dec. 8, 2020 must provide a privacy label}. \emph{Old apps without a label on 2 April 2021}: since crawling began four months after the policy, 266,740 old apps (22.9\,\%) already had a label in the first snapshot, so the time of their first compliance is unknown; RQ1/RQ2 analyses mostly use the remaining 77.1\,\%. \emph{Ratios}: the weekly app count fluctuates around 1.2 million and the number of old apps falls, so most plots use weekly ratios with a stated, non-constant denominator (e.g.\ old apps without a label on 2 April) \pc{li2022}{3}{we chose to use a weekly ratio over the raw number to describe the trends in most of the subsequent plots}.

\cf{Results RQ1} 51.6\,\% of all apps still had no label on 5 Nov 2021, eleven months after enactment; new apps are 100\,\% compliant by construction, but only 35.3\,\% of old apps created a label and the weekly number of old apps adding a first label decreases \pc{li2022}{3}{Due to Apple’s requirement, all new apps published after Dec. 8, 2020 must have a privacy label to enter the App Store, resulting in an 100\% compliance rate}. Only 2.7\,\% of old apps created a label without a simultaneous update (``voluntary adoption'') \pc{li2022}{4}{only 2.7\% of old apps created a privacy label without simultaneous app updates}; of those that created a first label over the seven months, 64.2\,\% updated the app at the same time and 14.4\,\% within one week after \pc{li2022}{4}{14.4\% made an app update within one week}. Interpretation: little incentive for inactive apps.

\cf{Results RQ2} Restricted to the 137,088 apps that created their first label in April 2021, giving a fixed observation window of 28 weeks \pc{li2022}{4}{we only considered the 137,088 apps that created the}: only 7,884 (5.8\,\%) ever updated the label, but 59,555 (43.4\,\%) released at least one app update \pc{li2022}{4}{only 7,884 of them ever updated the privacy label later (5.8\%)}. The gap suggests labels are not kept in step with data-practice changes, though app updates do not always change practices.

\cf{Results RQ3} Among labelled apps the share reporting ``Data Not Collected'' is stable; among apps reporting collection, the three categories ``Data Not Linked to You'', ``Data Linked to You'' and ``Data Used to Track You'' are not mutually exclusive (an app counts for a category if it names at least one data type under it) \pc{li2022}{4}{Note that these three categories are not mutually exclusive and we count an app as reporting a certain type of data practice}. Tracking declined from late April 2021 when iOS 14.5 introduced App Tracking Transparency, which the authors read as a possibly positive effect (a temporal association, not a causal estimate). Breakdown by the 26 store categories shows very different patterns of data use \pc{li2022}{4}{There are 26 app categories in total on the App Store}. An illustrative text test: 982 apps whose release notes mention ``cloud storage'' were searched for mismatches, and one location-logging app that stores data in the cloud claimed ``Data Not Collected'' \pc{li2022}{7}{found 982 apps that contained the keyword cloud storage}.

\cf{Discussion and limits} Privacy-label update does not follow app update, and the store shows no further alert once a label exists \pc{li2022}{6}{Currently, the App Store does not show further alerts}. Authors' limits: the ``date published'' field may not equal the true first-release date (1.3\,\% of new apps have an earlier first-version date; 0.52\,\% of new apps lack a label), and the late start of crawling excludes 22.9\,\% of old apps from RQ1/RQ2 so trends may not generalise \pc{li2022}{7}{because we started crawling four months after the requirement of privacy labels}. Labels may contain errors, and systematic errors ``remain to be investigated'' \pc{li2022}{7}{Although we want to note that these labels may contain errors}. Card additions: single country (U.S.), single store (iOS), no ground truth, no test of whether labels match behaviour, no uncertainty estimates, and a short preliminary format.

\cf{Use for this thesis} (1) A model of how to describe a store-wide capture in stage D: snapshot cadence, error-check and re-crawl step, the time window, the count of apps, and the denominators used for each ratio. (2) A reminder to define cohorts that are affected differently by a regulation (old vs.\ new, labelled vs.\ unlabelled) before comparing apps. (3) Evidence that the App Store label is a self-report, often missing or stale, so the thesis must not treat label content as ground truth for the app's behaviour; label--screen--policy mismatches are themselves a finding to code. (4) The ``not mutually exclusive'' counting rule for label categories is a pattern to adopt for multi-label privacy codes. The paper offers no protocol for reproducing the dataset (the crawler is not described in detail), so the thesis' own capture protocol must be written out in full and dated, since labels change.
\end{papercard}
